\chapter{Central Limit, Kepler Orbits and Sommerfeld Holonomy}
\label{chap:limite-central-kepler-sommerfeld}
\index{Kepler, orbit of}\index{Sommerfeld, quantization of}
\index{Levi-Civita, covering of}\index{Möbius, holonomy of}

This chapter distinguishes four strata that physical terminology may improperly superimpose: the exact cancellation of interior flows in a graph; the continuous realization of a central field; the Keplerian dynamics resulting from the inverse potential; and the semiclassical phase condition on a band with holonomy. The first is combinatorial. The other three successively require a metric chart, a Lagrangian, and a semiclassical phase line.

\section{Discrete Gauss and inverse descent to area}
\label{sec:gauss-discreto-limite-central}

Let \(G=(V,E)\) be a finite undirected graph, let
\(\vec E:=\{(x,y),(y,x):\{x,y\}\in E\}\) be its set of arcs, and let
\(F:\vec E\to\mathbb R\) be an antisymmetric \(1\)-co\-chain,
\(F(y,x)=-F(x,y)\). We define
\[
 \operatorname{div}F(x)
 :=\sum_{y:\{x,y\}\in E}F(x,y).
\]
For \(S\subseteq V\), the flux
\(\operatorname{Flux}_{\partial S}F\) is the oriented sum over the edges
with one endpoint in \(S\) and the other in \(V\setminus S\).

\begin{theorem}[Gauss's Theorem in Graphs]
\label{thm:gauss-grafo-hmt-md}
For every \(S\subseteq V\),
\[
 \sum_{x\in S}\operatorname{div}F(x)
 =\operatorname{Flux}_{\partial S}F.
 \tag{K1}\label{eq:gauss-grafo-hmt-md}
\]
\end{theorem}

\begin{proof}
Every edge whose two endpoints lie in \(S\) occurs twice with opposite
orientations and cancels. Only the edges crossing the boundary chain
remain.
\end{proof}

In the cubic lattice, for \(R\in\mathbb Z_{\geq0}\),
\[
 C_R:=\{x\in\mathbb Z^3:\|x\|_\infty\leq R\},
\]
the number of outgoing links is
\[
 |\partial C_R|=6(2R+1)^2.
\]
If a source has total flux \(Q\) and mean uniformity is imposed on the
shell, the flux per link is
\[
 E_R=\frac{Q}{6(2R+1)^2}.
\]
When defining \(r_{\rm eff}\) by
\[
 4\pi r_{\rm eff}^2=6(2R+1)^2,
\]
the exact identity is obtained
\[
 E_R=\frac{Q}{4\pi r_{\rm eff}^2}.
 \tag{K2}\label{eq:inverso-area-cascaron}
\]
The passage from \eqref{eq:gauss-grafo-hmt-md} to
\eqref{eq:inverso-area-cascaron} uses the additional hypothesis of mean
isotropy. Identifying \(Q\) with a gravitational or electric source also requires
the corresponding dimensional transducer.

\section{Central Lagrangian and Keplerian conics}
\label{sec:lagrangiano-central-kepler}

Let \(\mu>0\) be a reduced mass and
\(\mathcal K>0\) a coupling of the energy-times-length type,
\([\mathcal K]=\mathsf E\,\mathsf L\). On the continuous
Euclidean sheet, consider
\[
 L_{\rm Kep}(\mathbf r,\dot{\mathbf r})
 =\frac{\mu}{2}\|\dot{\mathbf r}\|^2+\frac{\mathcal K}{r},
 \qquad
 H_{\rm Kep}(\mathbf r,\mathbf p)
 =\frac{\|\mathbf p\|^2}{2\mu}-\frac{\mathcal K}{r},
 \tag{K3}\label{eq:lagrangiano-hamiltoniano-kepler}
\]
where \(r=\|\mathbf r\|>0\) and
\(\mathbf p=\mu\dot{\mathbf r}\). The Euler--Lagrange equation is
\[
 \mu\ddot{\mathbf r}
 =-\frac{\mathcal K}{r^3}\mathbf r.
 \tag{K4}\label{eq:ecuacion-central-kepler}
\]
The specific angular momentum
\(\mathbf h=\mathbf r\times\dot{\mathbf r}\) is constant; if
\(\mathbf h\ne0\), the trajectory remains in the plane orthogonal to
\(\mathbf h\).

\begin{proposition}[Keplerian conic]
\label{prop:conica-kepleriana-hmt-md}
Let \(h=\|\mathbf h\|\). Every nonradial solution of
\eqref{eq:ecuacion-central-kepler} satisfies
\[
 r(\phi)
 =\frac{p}{1+e\cos(\phi-\phi_0)},
 \qquad
 p=\frac{\mu h^2}{\mathcal K}.
 \tag{K5}\label{eq:conica-kepleriana}
\]
For negative energy, \(0\leq e<1\) and the orbit is elliptic.
\end{proposition}

\begin{proof}
With \(u=1/r\), the Binet equation takes the form
\[
 u''(\phi)+u(\phi)=\frac{\mathcal K}{\mu h^2}.
\]
Its general solution is
\[
 u=\frac{\mathcal K}{\mu h^2}
 \bigl(1+e\cos(\phi-\phi_0)\bigr).
\]
Inverting \(u\) gives \eqref{eq:conica-kepleriana}; the binding condition
is equivalent to \(e<1\).
\end{proof}

The areolar velocity is
\[
 \frac{\mathrm dA}{\mathrm dt}=\frac h2.
 \tag{K6}\label{eq:segunda-ley-kepler}
\]
For an ellipse with semimajor axis \(a\), integrating
\eqref{eq:segunda-ley-kepler} over one period yields
\[
 A_{\rm ell}=\pi ab,
 \qquad
 b=a\sqrt{1-e^2},
 \qquad
 p=a(1-e^2)=\frac{\mu h^2}{\mathcal K}.
\]
therefore,
\(T=2A_{\rm ell}/h\), and to the eliminar \(b\), \(p\) and \(h\) is obtained
\[
 T^2=\frac{4\pi^2\mu}{\mathcal K}\,a^3.
 \tag{K7}\label{eq:tercera-ley-kepler}
\]
In the gravitational test problem,
\(\mathcal K=G M\mu\), and
\eqref{eq:tercera-ley-kepler} reduces to
\(T^2=4\pi^2a^3/(GM)\) \cite{Kepler1609,Newton1687}.

The graph Gauss theorem is exact. Under the declared isotropy and
distance, the flux produces the inverse-area field; the continuum limit
then transports that field to the central Lagrangian and to the
Keplerian orbits written above.

\section{Levi-Civita Covering and Möbius Strip}
\label{sec:levi-civita-mobius-kepler}

In the complex plane without the origin, the Levi--Civita application
\[
 q=z^2
 \tag{K8}\label{eq:cubierta-levi-civita}
\]
is a double cover. One turn of \(q\) around the origin interchanges
the two lifts \(z\) and \(-z\); two turns are required to close the
lift. It is
an exact witness of memory \(\mathbb Z_2\). It is not, however, a Möbius
band: full regularization of the Kepler problem also requires
reparameterizing time and transforming the momenta.

Now let \(\gamma:S^1\hookrightarrow\mathbb R^2\) be an embedding whose image is an elliptic orbit, and
fix \(\varepsilon>0\). We define
\[
 M_\gamma
 :=([0,2\pi]\times[-\varepsilon,\varepsilon])/
 \bigl((0,u)\sim(2\pi,-u)\bigr).
 \tag{K9}\label{eq:banda-mobius-orbita}
\]

\begin{theorem}[Orbital band oriented period]
\label{thm:periodo-orientado-banda-orbital}
\(M_\gamma\) is a Möbius band whose kernel is \(\gamma\). A section
of the sign line has holonomy
\[
 \mathcal H(\gamma)=-I,
 \qquad
 \mathcal H(\gamma^2)=I.
 \tag{K10}\label{eq:holonomia-banda-orbital}
\]
Therefore, its return as an oriented section requires two turns.
\end{theorem}

\begin{proof}
Identification of the extreme fibers by means of \(u\mapsto-u\) is the
standard construction of the Möbius band. The holonomy of one turn is
the fiber involution; its square is the identity.
\end{proof}

The cover \eqref{eq:cubierta-levi-civita} and the band
\eqref{eq:banda-mobius-orbita} share a monodromy of order two, but they are not
isomorphic objects. This distinction prevents a particle--antiparticle
orientation from being deduced from Levi--Civita regularization without an
additional map.

\section{Sommerfeld's semiclassical condition with holonomy}
\label{sec:sommerfeld-holonomia-banda}

Let
\[
 J_\gamma:=\oint_\gamma\mathbf p\cdot\mathrm d\mathbf q
\]
the action of a cycle, and suppose that semiclassical transport contributes the
phase \(\exp(\mathrm iJ_\gamma/\hbar_{\rm int})\). If the associated real line
has holonomy \(\varepsilon_\gamma\in\{+1,-1\}\), the closure condition of
the complexified section is
\[
 \varepsilon_\gamma
 \exp\!\left(\frac{\mathrm iJ_\gamma}{\hbar_{\rm int}}\right)=1.
 \tag{K11}\label{eq:cierre-fase-holonomia-sommerfeld}
\]

\begin{proposition}[Shift by half a unit due to holonomy of sign]
\label{prop:desplazamiento-sommerfeld-holonomia}
Under the preceding hypotheses,
\[
 \begin{array}{lll}
 \varepsilon_\gamma=+1
 &\Longrightarrow&
 J_\gamma=2\pi\hbar_{\rm int}n,\\[1mm]
 \varepsilon_\gamma=-1
 &\Longrightarrow&
 J_\gamma=2\pi\hbar_{\rm int}\left(n+\dfrac12\right),
 \end{array}
 \qquad n\in\mathbb Z.
 \tag{K12}\label{eq:cuantizacion-sommerfeld-holonomia}
\]
\end{proposition}

\begin{proof}
The \cref{eq:cierre-fase-holonomia-sommerfeld} requires a phase equal to one in the trivial sector and
equal to minus one in the sign sector. Their arguments differ,
respectively, by an integral or half-integral multiple of
\(2\pi\).
\end{proof}

The proposition is an exact consequence of the declared phase and
holonomy; it does not by itself select the Coulomb Hamiltonian, a physical
quantum number, or an atomic spectrum \cite{Sommerfeld1916}. In an
EBK quantization, the Maslov indices of the corresponding cycles must be
added without counting the topological shift twice.

Relative to the orientation reader, the two local sheets may
represent routes \(\gamma\) and \(\gamma^\vee\) and be linked to the
two orientations of the two-sheet entanglement defined in
\cref{def:entropia-bicapa-rev6}. This correspondence neither identifies
those routes with advanced or retarded Green functions nor asserts that
entropy causes the Keplerian conic: the orbit proceeds from
\eqref{eq:ecuacion-central-kepler}; entropy measures, after a
Hilbert completion, correlation between sheets.

\begin{criterion}[Kepler--Sommerfeld limit]
\label{crit:limite-kepler-sommerfeld}
The realization is refuted or improperly typed if an inverse-square law is obtained
without declaring isotropy; if
\(\mathcal K\) or \(G\) is introduced as a scalar without a dimensional line;
if cover \(q=z^2\) is identified with a Mobius band; if
\eqref{eq:cuantizacion-sommerfeld-holonomia} is published as an exact spectrum
without a Hamiltonian, cycles, and Maslov indices; or if
entanglement is attributed to a merely deterministic route inversion.
\end{criterion}
